\subsection{Qué hacía el TP x86}

El TP3 usaba \textbf{task switching por hardware}: un descriptor de TSS en la GDT por tarea
(\texttt{tss\_initGdtEntry}), una TSS por tarea (\texttt{fill\_tss\_idle}, \texttt{fill\_tss}) con \texttt{eip},
\texttt{cr3}, \texttt{esp}, \texttt{eflags}, selectores de segmento, etc. El cambio de tarea se hacía con un \texttt{jmp
<sel>:<off>} al selector de la TSS, dejando que el procesador guardara/cargara todo el contexto automáticamente.

\subsection{Equivalente en RISC-V}

RISC-V no tiene task switching por hardware, todo es por software. Se simplifica bastante: no hay GDT, ni TSS, ni
descriptores; queda una estructura por tarea y nada más:

\begin{codesnippet}
\begin{verbatim}
 struct Task {
   TrapFrame tf;
   uint32_t initialized;
 };
 static Task g_tasks[20];
 static TrapFrame g_idle_tf;
\end{verbatim}
\end{codesnippet}

Cada \texttt{TrapFrame} cumple el rol que en el TP x86 cumplía la TSS: contiene el \emph{snapshot} completo del estado
de esa tarea. El campo \texttt{initialized} replica la noción de ``descriptor presente''.

\subsection{Inicialización de \texttt{TrapFrame} (equivalente de \texttt{fill\_tss})}

Toda la ``configuración inicial de la TSS'' del TP x86 (eip, esp, cs/ds/ss, eflags) se colapsa en tres campos del
\texttt{TrapFrame}:

\begin{codesnippet}
\begin{verbatim}
 static void init_user_trapframe(TrapFrame& tf, uint64_t entry_pc) {
     tf = {};
     tf.sepc = entry_pc;              // PC inicial
     tf.sp   = TASK_STACK;            // stack de usuario
     // SPIE=1 (SIE <- SPIE tras sret), SPP=0 (retorno a U-mode).
     tf.sstatus = (1ull << 5);
 }
\end{verbatim}
\end{codesnippet}

\begin{itemize}
  \item \texttt{sepc}$=$\texttt{TASK\_CODE} (\texttt{0x08000000}) $\Leftrightarrow$ \texttt{tss.eip}$=$dirección inicial
  del código.
  \item \texttt{sp}$=$\texttt{TASK\_STACK}$=$\texttt{TASK\_CODE}+2$\cdot$4~KiB $\Leftrightarrow$ \texttt{tss.esp}.
  \item \texttt{sstatus.SPIE}$=$1 asegura que \texttt{SIE}$=$1 luego del \texttt{sret}, es decir, la tarea corre con
  interrupciones habilitadas (equivalente a \texttt{eflags.IF}$=$1 del TP x86, donde se usaba \texttt{0x202}).
  \item \texttt{sstatus.SPP}$=$0 indica que el \texttt{sret} devuelve el control a U-mode (equivalente a ``DPL=3''). No
  hay ``CS'' ni ``DS''; el modo es el único eje de privilegio.
\end{itemize}

Nótese que en RISC-V no existen ``registros de segmento'': eso elimina por completo los campos
\texttt{cs/ds/ss/fs/gs/es} de la TSS. Tampoco hace falta una ``stack de nivel 0'' por tarea (\texttt{ss0:esp0}): al
entrar al kernel vía trap, se cambia a el stack de trap única (\texttt{trap\_stack\_top}), que alcanza porque SIE se
limpia al entrar al trap (no hay traps anidados al mismo nivel de privilegio).

\subsection{``Salto a la tarea Idle''}

El TP x86 arrancaba con \texttt{jmp 0xB8:0} después de cargar \texttt{tr} con el selector de la tarea inicial. En este
port, \texttt{kernel\_main} hace:

\begin{codesnippet}
\begin{verbatim}
 init_idle_once();

 asm volatile("csrw stvec, %0" : : "r"((uint64_t)trap_entry));
 const uint64_t sie_mask = (1ull << 1) | (1ull << 5); // SSIE + STIE
 asm volatile("csrs sie, %0" : : "r"(sie_mask));

 TrapFrame* first = schedule_next();
 enter_user(first);
\end{verbatim}
\end{codesnippet}

\texttt{enter\_user} (en \texttt{trap.S}) restaura todos los GPRs del \texttt{TrapFrame} y ejecuta \texttt{sret}. Con
\texttt{sstatus.SPP}$=$0 y \texttt{sepc}$=$\texttt{TASK\_CODE}, esto transiciona a U-mode saltando al primer byte del
binario de idle.

\subsection{La tarea Idle}

En el TP x86, idle era un binario del kernel en \texttt{0x00014000} que compartía \texttt{cr3} con el kernel. En este
port idle es un binario de usuario \emph{separado} (\texttt{riscv/user/idle.c}, ver sección~\ref{sec:user}), con su
propia address space (\texttt{asid}$=$0) y su propio buffer de 8~KiB como ``celda'' (no está en el tablero). Su código
completo es:

\begin{codesnippet}
\begin{verbatim}
 void user_main() {
     for (;;) __asm__ volatile("nop");
 }
\end{verbatim}
\end{codesnippet}

El porqué de hacerla U-mode y no S-mode: mantiene el mismo ABI/path de entrada/salida que los mineros, simplifica el
\texttt{TrapFrame} y evita tener que distinguir ``contexto supervisor'' vs ``contexto usuario'' en el scheduler.
